\scanpage{71}
% archive foliation: 70
\uline{Digression} Épinglages de revêtements par \deleted{\unsure{vecteurs} fibrés tangents} rev. kummériens des \unsure{Espaces|espaces} tangents (cas relatif)

\rule{\linewidth}{0.4pt}

Soient
\[
\tag{1}f:X\longrightarrow S
\]
un morphisme lisse de dim relative $1$, $t$ une section de $X$ sur $S$, $X^*=X\setminus t(S)$.

On considère des schémas \inserted{$X'$} finis sur $X$, lisses sur $S$ de dim relative $1$, munis d'une section $t'$ au dessus de $t$, et ``modérément ramifiés'' en $t'$ sur $X$, avec un degré de ramification $n$ -- ce qui signifie que loc pour la top. étale sur $X,X'$, $X'$ est \unsure{isom} au voisinage de tt point $x'$ de $t'(S)$ à un $X$-schéma Kummérien de $X$ au voisinage de $x=p(x')$, de la forme
\margin{\uline{NB} $n$ est premier aux car. résiduelles de $S$ !!!}
\margin{$p:X'\longrightarrow X$}
\[
X'_0=\operatorname{Spec}\underline O_X[z]/(z^n-\varphi)
\]
où $\varphi$ est une équation locale du diviseur relatif $t(S)\subset X$ au voisinage de $x$.

Sous ces conditions, via l'hom norme locale, on définit un isom. canonique
\[
\tag{2}T_{X'/\unsure{S},t'}^{\otimes n}\xrightarrow[\sim]{\kappa_{t'}}T_{X/S,t}
\]
